% Independent audit appendix; not a claim that the missing instance is proved.
\section{The source-to-command interface}
\label{sec:source-command-audit}

The literal physical certificate and the abstract command compiler have
different endpoints.  The former simulates one resolved increment,
decrement, jump, or halt instruction of the nine-register source.  The
latter starts from the already specified arithmetic meanings of
$\mathrm{AX}_{c,x,y,z}$, $\mathrm{MP}_x$, and $S$.  A complete connection
between them needs the following intermediate lemma.

Let $\operatorname{SrcStep}$ be the totalized step function of the exact
3,671-instruction source.  For a live scheduler configuration
$c=(n,\ell,t,w,d,h)$ define
\[
 J(c)=\bigl(25;(n,t,w,0,d,h,0,0,\ell)\bigr).
\]
Thus the auxiliary registers $A,X,Y$ are zero at the common fetch port.
The required claim is that there is an explicitly specified primitive
recursive $\theta$, positive on these configurations, such that
\[
 \mathrm{iPRA}\vdash
 \operatorname{SrcStep}^{\theta(c)}(J(c))=J(R_{11H}(c)).
 \tag{SC}
\]
Here either the abstract scheduler's halted configurations and $J$ must
retain enough information to describe the exact terminal source registers,
or both sides must explicitly collapse every halted configuration to a
single absorbing halt code, as in the physical bridge below. Merely
adding an absorbing flag while retaining unspecified terminal registers
would not specify an exact endpoint. The initial
configuration is $J(2,0,0,0,0,0)$.  The Q variant changes only the finite
AX expression.  Once (SC) and its halt-flag agreement are supplied, the
positive simulation theorem applies again, now between the command and
primitive-source levels.

\subsection{The inherited arithmetic is sufficient in principle}

The necessary loop bounds are concrete, not speculative unbounded
termination arguments.  For example, the inherited Cantor decoder uses
\[
 L(0,p)=1,\qquad L(m+1,p)=r(p)+L(m,\operatorname{advance}(p)),
\]
where $r(a,b)=5$ if $a>0$ and $r(0,b)=3b+8$.  Its two cases enumerate
the Cantor pairs.  On input $w$, the loop takes at most
$w(3w+8)+1$ steps; the final transfer takes $3\operatorname{tl}(w)+1$.
These are the functions and bounds used in the inherited
\path{M6/Source/Unpair.lean} and the retained constructor proof.
The pairing phase has the exact cost
\[
 P(t,a)=4t+2+\sum_{j<a+t}(7j+5),
\]
and the fourfold multiplier has cost $6y+1$.  A copy followed by restoration
has cost $7b+2$.  Clearing a register and the two comparison cleanup
loops also have linear costs.  Their endpoint equations follow from
quantifier-free counter invariants, with all unrelated registers carried
as parameters.  No induction over PA formulas is involved.

The inherited fetch lemma even supplies the bounds
\[
 2+\operatorname{decodeCost}(\ell-1)\quad(\ell>0),\qquad
 7n+4+\operatorname{decodeCost}(n)\quad(\ell=0),
\]
before the finite dispatch.  It proves exactly the exceptional fresh
candidate convention used in the new abstract proof.  Its cost functions
are primitive recursive.

Some inherited statements existentially quantify their duration, but the
proofs construct it by these linear costs and primitive recursions.  For
example, a simultaneous destructive comparison can use the recurrence
\[
 C(0,0)=2,\quad C(0,x+1)=3x+5,\quad
 C(h+1,0)=2h+6,\quad C(h+1,x+1)=C(h,x)+3.
\]
Equivalently, this is a piecewise linear function obtained by splitting
on $h=x$, $h<x$, and $h>x$.  Thus existential notation in the Lean
statement is not a mathematical obstacle to a primitive recursive clock.
The finite AX word composes these clocks at the successive stores.

\subsection{What the present certificates still have to bind}

The unchanged macro arithmetic is reusable after its actual source ports
and complete outgoing control edges have been identified.  The old
\path{PA196/Source/CommandTransfer.lean} proves such a statement for
another source: it has a 201-token AX body, three AX fields, other literal
addresses, and the old Shift/MP behavior.  It cannot instantiate (SC)
unchanged.

The new \path{check_commands.py} checks the 190-token symbolic AX
expression, selected node shapes, four-field parser identities, and
abstract command words.  It assigns arithmetic operations to macro names;
in particular, its guarded-cleanup comment assumes the corresponding
macro meaning.  Its successful execution alone therefore does not prove
the source realization of those meanings.

The remaining instance is precise and finite at its binding layer:
record the entry, exit and internal control edges of each called macro
in the new resolved source; attach its generic arithmetic endpoint and
clock; compose them through fetch, dispatch and the command.  The changed
cases requiring explicit treatment are the extra AX push followed by the
MP field reader, the no-push successful MP tail, and the cached-only Shift
path.  For the last case, clearing $T$, copying and restoring $H$, then
incrementing $X$ reaches $(T,H,X)=(h,h,1)$.  The old destructive equality
guard then halts exactly when $h=1$; otherwise it clears its operands and
the unchanged stack decoder completes the specified Shift.  For AX,
the additional MP field must consume $z$ and return $D=z$, including its
$z=0$ equality branch.  All unused registers and the exact residual
continuation must be included in the endpoint, not just the visible top.

These observations supply a direct route to a written mathematical
$\mathrm{iPRA}$ proof of (SC); they do not reveal a stronger-induction
obstruction.  Nevertheless, until the new macro-port bindings and composed
clock are supplied, the responsible statement is that the source-to-command
instance is an outstanding proof obligation.  A primitive-instruction
physical certificate and a free-algebra command certificate do not, by
themselves, discharge it.  This obligation is distinct from, and much more
local than, the still-open PA-refutation-to-numerical-calculus calibration.
